\documentclass[conference]{IEEEtran}
\IEEEoverridecommandlockouts
% The preceding line is only needed to identify funding in the first footnote. If that is unneeded, please comment it out.
\usepackage{cite}
\usepackage{amsmath,amssymb,amsfonts}
\usepackage{algorithmic}
\usepackage{graphicx}
\usepackage{textcomp}
\usepackage{xcolor}
\def\BibTeX{{\rm B\kern-.05em{\sc i\kern-.025em b}\kern-.08em
    T\kern-.1667em\lower.7ex\hbox{E}\kern-.125emX}}
\begin{document}

\title{Code Reviewer AI}

\author{\IEEEauthorblockN{Aktuelle Themen und Trends I und II \\
BWSE - SS 2026}
\IEEEauthorblockA{\textit{Group Nb. 1} \\
Lara Kammerer 2310859024\\
Julia Michler 2510859031\\
Stefan Zuza 2410859032}
}

\maketitle

\begin{abstract}
Tick the boxes for the level your implementation achieves:

\begin{itemize}
    \item[$\square$] Minimal
    \item[$\boxtimes$] Advanced
    \item[$\square$] Expert
\end{itemize}

\end{abstract}

\section{Work Distribution}
All three team members jointly developed the static\_analyst during a dedicated project meeting, where the team also worked on the first attempts to integrate an MCP server into the system. In a separate follow-up meeting, the team collaboratively implemented the MCP server and, as part of this work, also created the repo\_setup\_agent to handle deterministic repository preparation, including repository resolution, cloning, structure inspection, and manifest summarization.

Lara contributed primarily to the implementation of the architecture\_design\_analyst and its supporting tools. In addition, she carried out preparatory work for the MCP integration, helped shape the technical setup of the server, and contributed to parts of the project documentation.

Julia was responsible for the implementation of the code\_quality\_documentation\_agent and its related tooling. Her work focused on readability, maintainability, documentation quality, and refactoring-oriented review support within the project.

Stefan implemented the security\_reviewer and its associated tools, as well as the performance\_optimizer. His responsibilities covered the design of security-focused analysis functionality and performance-oriented review capabilities for detecting vulnerabilities, risky coding practices, and inefficient implementation patterns.

\section{Problem Statement}
Our multi-agent system is designed to automate the analysis of software repositories by assigning different review tasks to specialized agents. The scenario is that a repository is provided through a `projects.json` file, after which the system prepares the repository and performs focused analyses such as static code review, architecture assessment, and, in extended versions, security, performance, documentation, and test coverage evaluation. The main goal is to make repository reviews more structured, reusable, and consistent.

A key assumption is that the repository can be resolved and cloned from the given input path and that all later analyses should work on the same prepared project state. We also assume that separating responsibilities across multiple specialized agents leads to clearer and more reliable results than using one single general-purpose agent.

\section{Architecture}
Describe the overall architecture of your multi-agent system.
Figures and diagrams illustrating the architecture or interaction flow are encouraged!

\subsection{Agent Description}
Our system is built around specialized agents that each focus on a specific part of the repository review process.

\begin{itemize}
	\item The \textbf{repo\_setup\_agent} is responsible for preparing the repository context. It resolves the repository from the input file, clones it, inspects the structure, and summarizes important manifest files. This agent acts as the entry point of the workflow and provides the shared project context for the other agents.
	\item The \textbf{static\_analyst} is responsible for deterministic static code analysis. Its task is to identify maintainability, style, and structural issues and to support its findings with concrete evidence from files and directories. It depends on the prepared repository context and uses this shared setup as the basis for its review.
	\item The \textbf{architecture\_design\_analyst} focuses on higher-level software design questions. It evaluates modularity, separation of concerns, coupling, dependency direction, and maintainability risks. Like the static analyst, it builds on the prepared repository context and uses the structural information gathered in the setup phase to assess the system architecture.
	\item \textbf{TODO}
\end{itemize}

\subsection{Task Description}
The system solves several repository review tasks that are distributed across specialized agents. 

\begin{itemize}
	\item First, the repository must be prepared: the input path is resolved, the repository is cloned, its structure is inspected, and important manifest files are summarized. This task is handled by the repo\_setup\_agent, which provides the shared context for the rest of the system.
	\item After the setup phase, the analysis tasks are divided by review type. The static\_analyst handles static code analysis and is responsible for identifying maintainability, style, and structural issues in the codebase. The architecture\_design\_analyst focuses on higher-level design questions such as modularity, separation of concerns, coupling, boundaries, and long-term maintainability. \textbf{TODO}
\end{itemize}

\subsection{Flow}
Explain the interaction flow of the system. Describe how agents communicate, 
coordinate, and complete tasks. Diagrams are welcome.

\section{Minimal}
Describe the minimal functionality of your system and how it satisfies the minimal requirements.

\section{Advanced}
Describe the additional features implemented for the advanced level and how they extend the minimal system.

\section{Expert}
Describe the expert-level features of your system and explain why they go beyond the advanced level.

\section{Reflection}
Reflect on the development of your system. Discuss challenges, design decisions, 
limitations, and possible improvements for future work.

\section{Feedback - Optional}
You are invited to provide optional feedback on the course and the project.
Your comments will help improve the course for future students. Feel free to share what you found helpful, what could be improved, and any suggestions you may have. Thank you!



\subsection{References}

Please number citations consecutively within brackets \cite{Vaswani2017}. The 
sentence punctuation follows the bracket \cite{Vaswani2017}. Refer simply to the reference 
number, as in \cite{Vaswani2017}---do not use ``Ref. \cite{Vaswani2017}'' or ``reference \cite{Vaswani2017}'' except at 
the beginning of a sentence: ``Reference \cite{Vaswani2017} was the first $\ldots$''

%-----------------------------------
\bibliographystyle{IEEEtran}
\bibliography{library.bib}


\end{document}
